% =============================================================================
% Appendix: the anchor rule and the anchor codes (\label{sec:wm-app-anchors}).
% GENERATED by gonogo/phase05/make_anchor_appendix.py from the frozen
% outputs/phase05/freeze/label_sets.json (frozen 2026-11-03, SHA-256
% c1feffa3508d9070af4994544c08d62efcc78d0a1e48732e1716b3816cf4bce4). Do not edit by hand: edit the generator and rerun it.
% Meant to be included with % =============================================================================
% Appendix: the anchor rule and the anchor codes (\label{sec:wm-app-anchors}).
% GENERATED by gonogo/phase05/make_anchor_appendix.py from the frozen
% outputs/phase05/freeze/label_sets.json (frozen 2026-11-03, SHA-256
% c1feffa3508d9070af4994544c08d62efcc78d0a1e48732e1716b3816cf4bce4). Do not edit by hand: edit the generator and rerun it.
% Meant to be included with % =============================================================================
% Appendix: the anchor rule and the anchor codes (\label{sec:wm-app-anchors}).
% GENERATED by gonogo/phase05/make_anchor_appendix.py from the frozen
% outputs/phase05/freeze/label_sets.json (frozen 2026-11-03, SHA-256
% c1feffa3508d9070af4994544c08d62efcc78d0a1e48732e1716b3816cf4bce4). Do not edit by hand: edit the generator and rerun it.
% Meant to be included with % =============================================================================
% Appendix: the anchor rule and the anchor codes (\label{sec:wm-app-anchors}).
% GENERATED by gonogo/phase05/make_anchor_appendix.py from the frozen
% outputs/phase05/freeze/label_sets.json (frozen 2026-11-03, SHA-256
% c1feffa3508d9070af4994544c08d62efcc78d0a1e48732e1716b3816cf4bce4). Do not edit by hand: edit the generator and rerun it.
% Meant to be included with \input{method_appendix_anchors} in the appendix of
% the main file; labels use the wm namespace (sec:wm-app-*, tab:wm-app-*).
% =============================================================================

\section{The anchor rule and the anchor codes}
\label{sec:wm-app-anchors}

The anchor set $\mathcal{C}_0$ of \S\ref{sec:wm-deconf} is produced by a rule that uses label statistics alone and that is frozen, together with the codes it selects, before the pre-specified analyses. Its form, a broad rule-based set rather than a short curated list, was informed by exploratory analyses of the same data. A code is a candidate when it is a non-cardiac three-character diagnostic code with at least 200 transitions, onsets and resolutions together, on the pairs of the training folds. Transitions in both directions count, because the context estimator receives both as indicators, and because the set works by breadth: individual anchors are weak proxies of the context, which is recovered from many of them jointly.

A candidate is excluded when it has a plausible route to the electrocardiogram, which we read broadly. The excluded groups are electrolyte, fluid, nutritional and renal disorders; endocrine and metabolic conditions with electrocardiographic effects, and cardiovascular risk factors; conditions that change the heart rate through fever, sepsis, anaemia, bleeding or hypoxia; pulmonary disease and sleep apnoea; drugs with electrocardiographic effects, chemotherapy and transplantation; codes with cardiac content, such as cardiac symptoms, findings, devices, history and procedures, and codes linked to syncope; disorders of the brain, stroke, seizures and neuromuscular disease; liver disease; gastroparesis and intestinal ischaemia; systemic inflammatory disease; malignancy; pregnancy; and psychiatric and substance-use disorders, through psychotropic drugs, stimulants and alcohol. Codes whose strings are ambiguous between the coding systems from which the labels were converted are excluded as well. Every other candidate is a no-route code.

Each no-route code is then screened for separation, the second requirement of \S\ref{sec:wm-deconf}. We compute the Mantel--Haenszel odds ratio of an onset of any grounded cardiac code, namely the positive controls, myocardial infarction and pulmonary embolism, on pairs on which the candidate has an onset against pairs on which it has none, stratified by the care setting of the later encounter, by age tertile and by sex. A candidate whose odds ratio exceeds 1.5 travels with grounded events beyond what the encounter explains; it leaves $\mathcal{C}_0$ and is treated as a code with a route (Table~\ref{tab:wm-app-screen}). The screen is applied to every candidate. Its odds ratio is imprecise for rare codes, so some candidates may fail by chance, which only removes them.

The codes that pass form $\mathcal{C}_0$ (Table~\ref{tab:wm-app-anchors}). Each is assigned a recording profile. A code is recorded in every care setting when at least 25 of its onsets on pairs without a cardiac change fall into encounters summarized only by an emergency discharge, and it is recorded only during admissions otherwise. Codes of the first profile mostly mark acute complaints and track the acute context, while codes of the second track coding intensity and provenance; the profile also decides to which secondary margin a code contributes (\S\ref{sec:wm-theory}). A fixed subset of $\mathcal{C}_0$, drawn by an earlier seeded split, serves only as negative controls that validate the procedure, and it never enters the context estimator. These codes were examined in an earlier analysis of the same data, so validation on them is not blind. The remaining codes are rotated. They are split into five groups, stratified by their number of onsets on pairs without a cardiac change: the codes, sorted by that number, are cut into consecutive blocks of five, and the five group labels are permuted within each block with a fixed seed (Table~\ref{tab:wm-app-groups}). In each rotation one group is withheld from the context estimator and given code channels, and the other groups enter it. The rule, the reason recorded for every code, the screen results and the groups are stored in one file whose checksum was fixed before any statistic that depends on them was computed.

\begin{table}[t]
\centering
\small
\caption{Anchor codes $\mathcal{C}_0$ selected by the rule of Appendix~\ref{sec:wm-app-anchors}, by ICD-10 chapter. Bold: the validation subset, used only as negative controls. $^{\ast}$: recorded in every care setting; all other codes are recorded only during admissions.}
\label{tab:wm-app-anchors}
\begin{tabular}{p{0.27\linewidth}p{0.66\linewidth}}
\hline
\textbf{Chapter} & \textbf{Codes} \\
\hline
Infectious and parasitic diseases & B25, B34$^{\ast}$, B35, B37, B49, B95, B96, B97$^{\ast}$ \\
Neoplasms & D12, D25 \\
Blood and immune mechanism & D73, D80, D89 \\
Endocrine, nutritional and metabolic & E04, E29, E53, E55, E74, E80 \\
Nervous system & \textbf{G43$^{\ast}$}, G44, G51, G56, G57, G58, G60, G89$^{\ast}$ \\
Eye and adnexa & H02, H04, H10, H26, H49, H53$^{\ast}$, H54, H57 \\
Ear and mastoid & H91, H93 \\
Respiratory system & J01, J02$^{\ast}$, \textbf{J06$^{\ast}$}, J20$^{\ast}$, J30, J32, J34, J40$^{\ast}$ \\
Digestive system & K04, K08, K12, K13, K20, \textbf{K21$^{\ast}$}, \textbf{K22}, K25, K26, K27, \textbf{K29$^{\ast}$}, K30, K40, K42, K43, K44, K50, K51, \textbf{K57}, K58, \textbf{K59$^{\ast}$}, K62$^{\ast}$, K63, \textbf{K64}, K66, K68, \textbf{K80}, K81, K82, K83, K91, K94 \\
Skin and subcutaneous tissue & L02, L03$^{\ast}$, L25, L27, L29, L30, L40, L50, L53, L59, L73, L89, L98 \\
Musculoskeletal system & M00, M11, M12, M14, M15, M16, M17, M19, M21, M25$^{\ast}$, M40, M41, M43, M46, M47, M48, M50, M51, M54$^{\ast}$, M75, M79$^{\ast}$, M81, M84, M85, M86, M87, M90 \\
Genitourinary system & N10, N12, N20, N30, N31, N32, N39$^{\ast}$, \textbf{N40}, N50, N52, N83, N92 \\
Symptoms, signs and abnormal findings & R10$^{\ast}$, R11$^{\ast}$, R14, R15, R21, R22, R30, R32, R33, R39, R49, R51$^{\ast}$, R52, R59, R61, R76, R82 \\
Injury, poisoning and other external consequences & S10, S12, S20$^{\ast}$, S22, S30, S32, S40, S42, S52, S60, S62, S70, S72, S80, S82, S90, S91, S92, T36, T37, T39, T78$^{\ast}$, T80, T83, T84, T85 \\
External causes & X58$^{\ast}$, Y92$^{\ast}$, Y93, Y99 \\
Factors influencing health status and contact with health services & Z16, Z22, Z23, Z28, Z43, Z48, Z59, Z63, Z74, Z76, Z77, Z80, Z81, Z88, Z90, Z93, \textbf{Z96} \\
\hline
\end{tabular}
\end{table}

\begin{table}[t]
\centering
\small
\caption{No-route candidates removed by the separation screen of Appendix~\ref{sec:wm-app-anchors} and treated as codes with a route.}
\label{tab:wm-app-screen}
\begin{tabular}{p{0.93\linewidth}}
\hline
B00, B02, D18, D32, H90, J38, J39, K02, L76, N99, S37, T17, Z00, Z53, Z78 \\
\hline
\end{tabular}
\end{table}

\begin{table}[t]
\centering
\small
\caption{Rotation groups of the anchor codes outside the validation subset, stratified by the number of onsets on pairs without a cardiac change. In each rotation one group is withheld from the context estimator and given code channels.}
\label{tab:wm-app-groups}
\begin{tabular}{p{0.08\linewidth}p{0.85\linewidth}}
\hline
\textbf{Group} & \textbf{Codes} \\
\hline
1 & B37, B96, D73, E29, G57, G58, G89, H49, H91, K04, K08, K20, K27, K44, K83, L02, L30, L53, L59, M17, M21, M85, M86, N20, N83, R14, R21, R30, S20, S30, S40, S60, S82, T36, Y93, Z22, Z90 \\
2 & B97, E53, E80, G44, H04, H10, H26, J20, K50, K58, K94, L27, L89, M12, M14, M40, M41, M50, M79, N10, N32, N39, N50, N92, R52, S12, S32, S72, S80, S92, T37, T39, X58, Z16, Z23, Z77 \\
3 & B95, D12, E04, G51, H02, J30, J32, K13, K30, K43, K51, K62, K63, K82, L25, L50, L73, M51, M75, M81, M84, M90, N30, N31, R39, R49, R51, R61, R76, T78, T80, T85, Y92, Z43, Z74, Z80 \\
4 & B34, B35, B49, D25, D89, G56, H53, H57, J02, K25, K40, K42, K66, K68, K81, L29, L40, M00, M11, M19, M25, M43, M87, R11, R32, R33, R82, S10, S22, S42, S52, S62, S70, T84, Z48, Z59 \\
5 & B25, D80, E55, E74, G60, H54, H93, J01, J34, J40, K12, K26, K91, L03, L98, M15, M16, M46, M47, M48, M54, N12, N52, R10, R15, R22, R59, S90, S91, T83, Y99, Z28, Z63, Z76, Z81, Z88, Z93 \\
\hline
\end{tabular}
\end{table}
 in the appendix of
% the main file; labels use the wm namespace (sec:wm-app-*, tab:wm-app-*).
% =============================================================================

\section{The anchor rule and the anchor codes}
\label{sec:wm-app-anchors}

The anchor set $\mathcal{C}_0$ of \S\ref{sec:wm-deconf} is produced by a rule that uses label statistics alone and that is frozen, together with the codes it selects, before the pre-specified analyses. Its form, a broad rule-based set rather than a short curated list, was informed by exploratory analyses of the same data. A code is a candidate when it is a non-cardiac three-character diagnostic code with at least 200 transitions, onsets and resolutions together, on the pairs of the training folds. Transitions in both directions count, because the context estimator receives both as indicators, and because the set works by breadth: individual anchors are weak proxies of the context, which is recovered from many of them jointly.

A candidate is excluded when it has a plausible route to the electrocardiogram, which we read broadly. The excluded groups are electrolyte, fluid, nutritional and renal disorders; endocrine and metabolic conditions with electrocardiographic effects, and cardiovascular risk factors; conditions that change the heart rate through fever, sepsis, anaemia, bleeding or hypoxia; pulmonary disease and sleep apnoea; drugs with electrocardiographic effects, chemotherapy and transplantation; codes with cardiac content, such as cardiac symptoms, findings, devices, history and procedures, and codes linked to syncope; disorders of the brain, stroke, seizures and neuromuscular disease; liver disease; gastroparesis and intestinal ischaemia; systemic inflammatory disease; malignancy; pregnancy; and psychiatric and substance-use disorders, through psychotropic drugs, stimulants and alcohol. Codes whose strings are ambiguous between the coding systems from which the labels were converted are excluded as well. Every other candidate is a no-route code.

Each no-route code is then screened for separation, the second requirement of \S\ref{sec:wm-deconf}. We compute the Mantel--Haenszel odds ratio of an onset of any grounded cardiac code, namely the positive controls, myocardial infarction and pulmonary embolism, on pairs on which the candidate has an onset against pairs on which it has none, stratified by the care setting of the later encounter, by age tertile and by sex. A candidate whose odds ratio exceeds 1.5 travels with grounded events beyond what the encounter explains; it leaves $\mathcal{C}_0$ and is treated as a code with a route (Table~\ref{tab:wm-app-screen}). The screen is applied to every candidate. Its odds ratio is imprecise for rare codes, so some candidates may fail by chance, which only removes them.

The codes that pass form $\mathcal{C}_0$ (Table~\ref{tab:wm-app-anchors}). Each is assigned a recording profile. A code is recorded in every care setting when at least 25 of its onsets on pairs without a cardiac change fall into encounters summarized only by an emergency discharge, and it is recorded only during admissions otherwise. Codes of the first profile mostly mark acute complaints and track the acute context, while codes of the second track coding intensity and provenance; the profile also decides to which secondary margin a code contributes (\S\ref{sec:wm-theory}). A fixed subset of $\mathcal{C}_0$, drawn by an earlier seeded split, serves only as negative controls that validate the procedure, and it never enters the context estimator. These codes were examined in an earlier analysis of the same data, so validation on them is not blind. The remaining codes are rotated. They are split into five groups, stratified by their number of onsets on pairs without a cardiac change: the codes, sorted by that number, are cut into consecutive blocks of five, and the five group labels are permuted within each block with a fixed seed (Table~\ref{tab:wm-app-groups}). In each rotation one group is withheld from the context estimator and given code channels, and the other groups enter it. The rule, the reason recorded for every code, the screen results and the groups are stored in one file whose checksum was fixed before any statistic that depends on them was computed.

\begin{table}[t]
\centering
\small
\caption{Anchor codes $\mathcal{C}_0$ selected by the rule of Appendix~\ref{sec:wm-app-anchors}, by ICD-10 chapter. Bold: the validation subset, used only as negative controls. $^{\ast}$: recorded in every care setting; all other codes are recorded only during admissions.}
\label{tab:wm-app-anchors}
\begin{tabular}{p{0.27\linewidth}p{0.66\linewidth}}
\hline
\textbf{Chapter} & \textbf{Codes} \\
\hline
Infectious and parasitic diseases & B25, B34$^{\ast}$, B35, B37, B49, B95, B96, B97$^{\ast}$ \\
Neoplasms & D12, D25 \\
Blood and immune mechanism & D73, D80, D89 \\
Endocrine, nutritional and metabolic & E04, E29, E53, E55, E74, E80 \\
Nervous system & \textbf{G43$^{\ast}$}, G44, G51, G56, G57, G58, G60, G89$^{\ast}$ \\
Eye and adnexa & H02, H04, H10, H26, H49, H53$^{\ast}$, H54, H57 \\
Ear and mastoid & H91, H93 \\
Respiratory system & J01, J02$^{\ast}$, \textbf{J06$^{\ast}$}, J20$^{\ast}$, J30, J32, J34, J40$^{\ast}$ \\
Digestive system & K04, K08, K12, K13, K20, \textbf{K21$^{\ast}$}, \textbf{K22}, K25, K26, K27, \textbf{K29$^{\ast}$}, K30, K40, K42, K43, K44, K50, K51, \textbf{K57}, K58, \textbf{K59$^{\ast}$}, K62$^{\ast}$, K63, \textbf{K64}, K66, K68, \textbf{K80}, K81, K82, K83, K91, K94 \\
Skin and subcutaneous tissue & L02, L03$^{\ast}$, L25, L27, L29, L30, L40, L50, L53, L59, L73, L89, L98 \\
Musculoskeletal system & M00, M11, M12, M14, M15, M16, M17, M19, M21, M25$^{\ast}$, M40, M41, M43, M46, M47, M48, M50, M51, M54$^{\ast}$, M75, M79$^{\ast}$, M81, M84, M85, M86, M87, M90 \\
Genitourinary system & N10, N12, N20, N30, N31, N32, N39$^{\ast}$, \textbf{N40}, N50, N52, N83, N92 \\
Symptoms, signs and abnormal findings & R10$^{\ast}$, R11$^{\ast}$, R14, R15, R21, R22, R30, R32, R33, R39, R49, R51$^{\ast}$, R52, R59, R61, R76, R82 \\
Injury, poisoning and other external consequences & S10, S12, S20$^{\ast}$, S22, S30, S32, S40, S42, S52, S60, S62, S70, S72, S80, S82, S90, S91, S92, T36, T37, T39, T78$^{\ast}$, T80, T83, T84, T85 \\
External causes & X58$^{\ast}$, Y92$^{\ast}$, Y93, Y99 \\
Factors influencing health status and contact with health services & Z16, Z22, Z23, Z28, Z43, Z48, Z59, Z63, Z74, Z76, Z77, Z80, Z81, Z88, Z90, Z93, \textbf{Z96} \\
\hline
\end{tabular}
\end{table}

\begin{table}[t]
\centering
\small
\caption{No-route candidates removed by the separation screen of Appendix~\ref{sec:wm-app-anchors} and treated as codes with a route.}
\label{tab:wm-app-screen}
\begin{tabular}{p{0.93\linewidth}}
\hline
B00, B02, D18, D32, H90, J38, J39, K02, L76, N99, S37, T17, Z00, Z53, Z78 \\
\hline
\end{tabular}
\end{table}

\begin{table}[t]
\centering
\small
\caption{Rotation groups of the anchor codes outside the validation subset, stratified by the number of onsets on pairs without a cardiac change. In each rotation one group is withheld from the context estimator and given code channels.}
\label{tab:wm-app-groups}
\begin{tabular}{p{0.08\linewidth}p{0.85\linewidth}}
\hline
\textbf{Group} & \textbf{Codes} \\
\hline
1 & B37, B96, D73, E29, G57, G58, G89, H49, H91, K04, K08, K20, K27, K44, K83, L02, L30, L53, L59, M17, M21, M85, M86, N20, N83, R14, R21, R30, S20, S30, S40, S60, S82, T36, Y93, Z22, Z90 \\
2 & B97, E53, E80, G44, H04, H10, H26, J20, K50, K58, K94, L27, L89, M12, M14, M40, M41, M50, M79, N10, N32, N39, N50, N92, R52, S12, S32, S72, S80, S92, T37, T39, X58, Z16, Z23, Z77 \\
3 & B95, D12, E04, G51, H02, J30, J32, K13, K30, K43, K51, K62, K63, K82, L25, L50, L73, M51, M75, M81, M84, M90, N30, N31, R39, R49, R51, R61, R76, T78, T80, T85, Y92, Z43, Z74, Z80 \\
4 & B34, B35, B49, D25, D89, G56, H53, H57, J02, K25, K40, K42, K66, K68, K81, L29, L40, M00, M11, M19, M25, M43, M87, R11, R32, R33, R82, S10, S22, S42, S52, S62, S70, T84, Z48, Z59 \\
5 & B25, D80, E55, E74, G60, H54, H93, J01, J34, J40, K12, K26, K91, L03, L98, M15, M16, M46, M47, M48, M54, N12, N52, R10, R15, R22, R59, S90, S91, T83, Y99, Z28, Z63, Z76, Z81, Z88, Z93 \\
\hline
\end{tabular}
\end{table}
 in the appendix of
% the main file; labels use the wm namespace (sec:wm-app-*, tab:wm-app-*).
% =============================================================================

\section{The anchor rule and the anchor codes}
\label{sec:wm-app-anchors}

The anchor set $\mathcal{C}_0$ of \S\ref{sec:wm-deconf} is produced by a rule that uses label statistics alone and that is frozen, together with the codes it selects, before the pre-specified analyses. Its form, a broad rule-based set rather than a short curated list, was informed by exploratory analyses of the same data. A code is a candidate when it is a non-cardiac three-character diagnostic code with at least 200 transitions, onsets and resolutions together, on the pairs of the training folds. Transitions in both directions count, because the context estimator receives both as indicators, and because the set works by breadth: individual anchors are weak proxies of the context, which is recovered from many of them jointly.

A candidate is excluded when it has a plausible route to the electrocardiogram, which we read broadly. The excluded groups are electrolyte, fluid, nutritional and renal disorders; endocrine and metabolic conditions with electrocardiographic effects, and cardiovascular risk factors; conditions that change the heart rate through fever, sepsis, anaemia, bleeding or hypoxia; pulmonary disease and sleep apnoea; drugs with electrocardiographic effects, chemotherapy and transplantation; codes with cardiac content, such as cardiac symptoms, findings, devices, history and procedures, and codes linked to syncope; disorders of the brain, stroke, seizures and neuromuscular disease; liver disease; gastroparesis and intestinal ischaemia; systemic inflammatory disease; malignancy; pregnancy; and psychiatric and substance-use disorders, through psychotropic drugs, stimulants and alcohol. Codes whose strings are ambiguous between the coding systems from which the labels were converted are excluded as well. Every other candidate is a no-route code.

Each no-route code is then screened for separation, the second requirement of \S\ref{sec:wm-deconf}. We compute the Mantel--Haenszel odds ratio of an onset of any grounded cardiac code, namely the positive controls, myocardial infarction and pulmonary embolism, on pairs on which the candidate has an onset against pairs on which it has none, stratified by the care setting of the later encounter, by age tertile and by sex. A candidate whose odds ratio exceeds 1.5 travels with grounded events beyond what the encounter explains; it leaves $\mathcal{C}_0$ and is treated as a code with a route (Table~\ref{tab:wm-app-screen}). The screen is applied to every candidate. Its odds ratio is imprecise for rare codes, so some candidates may fail by chance, which only removes them.

The codes that pass form $\mathcal{C}_0$ (Table~\ref{tab:wm-app-anchors}). Each is assigned a recording profile. A code is recorded in every care setting when at least 25 of its onsets on pairs without a cardiac change fall into encounters summarized only by an emergency discharge, and it is recorded only during admissions otherwise. Codes of the first profile mostly mark acute complaints and track the acute context, while codes of the second track coding intensity and provenance; the profile also decides to which secondary margin a code contributes (\S\ref{sec:wm-theory}). A fixed subset of $\mathcal{C}_0$, drawn by an earlier seeded split, serves only as negative controls that validate the procedure, and it never enters the context estimator. These codes were examined in an earlier analysis of the same data, so validation on them is not blind. The remaining codes are rotated. They are split into five groups, stratified by their number of onsets on pairs without a cardiac change: the codes, sorted by that number, are cut into consecutive blocks of five, and the five group labels are permuted within each block with a fixed seed (Table~\ref{tab:wm-app-groups}). In each rotation one group is withheld from the context estimator and given code channels, and the other groups enter it. The rule, the reason recorded for every code, the screen results and the groups are stored in one file whose checksum was fixed before any statistic that depends on them was computed.

\begin{table}[t]
\centering
\small
\caption{Anchor codes $\mathcal{C}_0$ selected by the rule of Appendix~\ref{sec:wm-app-anchors}, by ICD-10 chapter. Bold: the validation subset, used only as negative controls. $^{\ast}$: recorded in every care setting; all other codes are recorded only during admissions.}
\label{tab:wm-app-anchors}
\begin{tabular}{p{0.27\linewidth}p{0.66\linewidth}}
\hline
\textbf{Chapter} & \textbf{Codes} \\
\hline
Infectious and parasitic diseases & B25, B34$^{\ast}$, B35, B37, B49, B95, B96, B97$^{\ast}$ \\
Neoplasms & D12, D25 \\
Blood and immune mechanism & D73, D80, D89 \\
Endocrine, nutritional and metabolic & E04, E29, E53, E55, E74, E80 \\
Nervous system & \textbf{G43$^{\ast}$}, G44, G51, G56, G57, G58, G60, G89$^{\ast}$ \\
Eye and adnexa & H02, H04, H10, H26, H49, H53$^{\ast}$, H54, H57 \\
Ear and mastoid & H91, H93 \\
Respiratory system & J01, J02$^{\ast}$, \textbf{J06$^{\ast}$}, J20$^{\ast}$, J30, J32, J34, J40$^{\ast}$ \\
Digestive system & K04, K08, K12, K13, K20, \textbf{K21$^{\ast}$}, \textbf{K22}, K25, K26, K27, \textbf{K29$^{\ast}$}, K30, K40, K42, K43, K44, K50, K51, \textbf{K57}, K58, \textbf{K59$^{\ast}$}, K62$^{\ast}$, K63, \textbf{K64}, K66, K68, \textbf{K80}, K81, K82, K83, K91, K94 \\
Skin and subcutaneous tissue & L02, L03$^{\ast}$, L25, L27, L29, L30, L40, L50, L53, L59, L73, L89, L98 \\
Musculoskeletal system & M00, M11, M12, M14, M15, M16, M17, M19, M21, M25$^{\ast}$, M40, M41, M43, M46, M47, M48, M50, M51, M54$^{\ast}$, M75, M79$^{\ast}$, M81, M84, M85, M86, M87, M90 \\
Genitourinary system & N10, N12, N20, N30, N31, N32, N39$^{\ast}$, \textbf{N40}, N50, N52, N83, N92 \\
Symptoms, signs and abnormal findings & R10$^{\ast}$, R11$^{\ast}$, R14, R15, R21, R22, R30, R32, R33, R39, R49, R51$^{\ast}$, R52, R59, R61, R76, R82 \\
Injury, poisoning and other external consequences & S10, S12, S20$^{\ast}$, S22, S30, S32, S40, S42, S52, S60, S62, S70, S72, S80, S82, S90, S91, S92, T36, T37, T39, T78$^{\ast}$, T80, T83, T84, T85 \\
External causes & X58$^{\ast}$, Y92$^{\ast}$, Y93, Y99 \\
Factors influencing health status and contact with health services & Z16, Z22, Z23, Z28, Z43, Z48, Z59, Z63, Z74, Z76, Z77, Z80, Z81, Z88, Z90, Z93, \textbf{Z96} \\
\hline
\end{tabular}
\end{table}

\begin{table}[t]
\centering
\small
\caption{No-route candidates removed by the separation screen of Appendix~\ref{sec:wm-app-anchors} and treated as codes with a route.}
\label{tab:wm-app-screen}
\begin{tabular}{p{0.93\linewidth}}
\hline
B00, B02, D18, D32, H90, J38, J39, K02, L76, N99, S37, T17, Z00, Z53, Z78 \\
\hline
\end{tabular}
\end{table}

\begin{table}[t]
\centering
\small
\caption{Rotation groups of the anchor codes outside the validation subset, stratified by the number of onsets on pairs without a cardiac change. In each rotation one group is withheld from the context estimator and given code channels.}
\label{tab:wm-app-groups}
\begin{tabular}{p{0.08\linewidth}p{0.85\linewidth}}
\hline
\textbf{Group} & \textbf{Codes} \\
\hline
1 & B37, B96, D73, E29, G57, G58, G89, H49, H91, K04, K08, K20, K27, K44, K83, L02, L30, L53, L59, M17, M21, M85, M86, N20, N83, R14, R21, R30, S20, S30, S40, S60, S82, T36, Y93, Z22, Z90 \\
2 & B97, E53, E80, G44, H04, H10, H26, J20, K50, K58, K94, L27, L89, M12, M14, M40, M41, M50, M79, N10, N32, N39, N50, N92, R52, S12, S32, S72, S80, S92, T37, T39, X58, Z16, Z23, Z77 \\
3 & B95, D12, E04, G51, H02, J30, J32, K13, K30, K43, K51, K62, K63, K82, L25, L50, L73, M51, M75, M81, M84, M90, N30, N31, R39, R49, R51, R61, R76, T78, T80, T85, Y92, Z43, Z74, Z80 \\
4 & B34, B35, B49, D25, D89, G56, H53, H57, J02, K25, K40, K42, K66, K68, K81, L29, L40, M00, M11, M19, M25, M43, M87, R11, R32, R33, R82, S10, S22, S42, S52, S62, S70, T84, Z48, Z59 \\
5 & B25, D80, E55, E74, G60, H54, H93, J01, J34, J40, K12, K26, K91, L03, L98, M15, M16, M46, M47, M48, M54, N12, N52, R10, R15, R22, R59, S90, S91, T83, Y99, Z28, Z63, Z76, Z81, Z88, Z93 \\
\hline
\end{tabular}
\end{table}
 in the appendix of
% the main file; labels use the wm namespace (sec:wm-app-*, tab:wm-app-*).
% =============================================================================

\section{The anchor rule and the anchor codes}
\label{sec:wm-app-anchors}

The anchor set $\mathcal{C}_0$ of \S\ref{sec:wm-deconf} is produced by a rule that uses label statistics alone and that is frozen, together with the codes it selects, before the pre-specified analyses. Its form, a broad rule-based set rather than a short curated list, was informed by exploratory analyses of the same data. A code is a candidate when it is a non-cardiac three-character diagnostic code with at least 200 transitions, onsets and resolutions together, on the pairs of the training folds. Transitions in both directions count, because the context estimator receives both as indicators, and because the set works by breadth: individual anchors are weak proxies of the context, which is recovered from many of them jointly.

A candidate is excluded when it has a plausible route to the electrocardiogram, which we read broadly. The excluded groups are electrolyte, fluid, nutritional and renal disorders; endocrine and metabolic conditions with electrocardiographic effects, and cardiovascular risk factors; conditions that change the heart rate through fever, sepsis, anaemia, bleeding or hypoxia; pulmonary disease and sleep apnoea; drugs with electrocardiographic effects, chemotherapy and transplantation; codes with cardiac content, such as cardiac symptoms, findings, devices, history and procedures, and codes linked to syncope; disorders of the brain, stroke, seizures and neuromuscular disease; liver disease; gastroparesis and intestinal ischaemia; systemic inflammatory disease; malignancy; pregnancy; and psychiatric and substance-use disorders, through psychotropic drugs, stimulants and alcohol. Codes whose strings are ambiguous between the coding systems from which the labels were converted are excluded as well. Every other candidate is a no-route code.

Each no-route code is then screened for separation, the second requirement of \S\ref{sec:wm-deconf}. We compute the Mantel--Haenszel odds ratio of an onset of any grounded cardiac code, namely the positive controls, myocardial infarction and pulmonary embolism, on pairs on which the candidate has an onset against pairs on which it has none, stratified by the care setting of the later encounter, by age tertile and by sex. A candidate whose odds ratio exceeds 1.5 travels with grounded events beyond what the encounter explains; it leaves $\mathcal{C}_0$ and is treated as a code with a route (Table~\ref{tab:wm-app-screen}). The screen is applied to every candidate. Its odds ratio is imprecise for rare codes, so some candidates may fail by chance, which only removes them.

The codes that pass form $\mathcal{C}_0$ (Table~\ref{tab:wm-app-anchors}). Each is assigned a recording profile. A code is recorded in every care setting when at least 25 of its onsets on pairs without a cardiac change fall into encounters summarized only by an emergency discharge, and it is recorded only during admissions otherwise. Codes of the first profile mostly mark acute complaints and track the acute context, while codes of the second track coding intensity and provenance; the profile also decides to which secondary margin a code contributes (\S\ref{sec:wm-theory}). A fixed subset of $\mathcal{C}_0$, drawn by an earlier seeded split, serves only as negative controls that validate the procedure, and it never enters the context estimator. These codes were examined in an earlier analysis of the same data, so validation on them is not blind. The remaining codes are rotated. They are split into five groups, stratified by their number of onsets on pairs without a cardiac change: the codes, sorted by that number, are cut into consecutive blocks of five, and the five group labels are permuted within each block with a fixed seed (Table~\ref{tab:wm-app-groups}). In each rotation one group is withheld from the context estimator and given code channels, and the other groups enter it. The rule, the reason recorded for every code, the screen results and the groups are stored in one file whose checksum was fixed before any statistic that depends on them was computed.

\begin{table}[t]
\centering
\small
\caption{Anchor codes $\mathcal{C}_0$ selected by the rule of Appendix~\ref{sec:wm-app-anchors}, by ICD-10 chapter. Bold: the validation subset, used only as negative controls. $^{\ast}$: recorded in every care setting; all other codes are recorded only during admissions.}
\label{tab:wm-app-anchors}
\begin{tabular}{p{0.27\linewidth}p{0.66\linewidth}}
\hline
\textbf{Chapter} & \textbf{Codes} \\
\hline
Infectious and parasitic diseases & B25, B34$^{\ast}$, B35, B37, B49, B95, B96, B97$^{\ast}$ \\
Neoplasms & D12, D25 \\
Blood and immune mechanism & D73, D80, D89 \\
Endocrine, nutritional and metabolic & E04, E29, E53, E55, E74, E80 \\
Nervous system & \textbf{G43$^{\ast}$}, G44, G51, G56, G57, G58, G60, G89$^{\ast}$ \\
Eye and adnexa & H02, H04, H10, H26, H49, H53$^{\ast}$, H54, H57 \\
Ear and mastoid & H91, H93 \\
Respiratory system & J01, J02$^{\ast}$, \textbf{J06$^{\ast}$}, J20$^{\ast}$, J30, J32, J34, J40$^{\ast}$ \\
Digestive system & K04, K08, K12, K13, K20, \textbf{K21$^{\ast}$}, \textbf{K22}, K25, K26, K27, \textbf{K29$^{\ast}$}, K30, K40, K42, K43, K44, K50, K51, \textbf{K57}, K58, \textbf{K59$^{\ast}$}, K62$^{\ast}$, K63, \textbf{K64}, K66, K68, \textbf{K80}, K81, K82, K83, K91, K94 \\
Skin and subcutaneous tissue & L02, L03$^{\ast}$, L25, L27, L29, L30, L40, L50, L53, L59, L73, L89, L98 \\
Musculoskeletal system & M00, M11, M12, M14, M15, M16, M17, M19, M21, M25$^{\ast}$, M40, M41, M43, M46, M47, M48, M50, M51, M54$^{\ast}$, M75, M79$^{\ast}$, M81, M84, M85, M86, M87, M90 \\
Genitourinary system & N10, N12, N20, N30, N31, N32, N39$^{\ast}$, \textbf{N40}, N50, N52, N83, N92 \\
Symptoms, signs and abnormal findings & R10$^{\ast}$, R11$^{\ast}$, R14, R15, R21, R22, R30, R32, R33, R39, R49, R51$^{\ast}$, R52, R59, R61, R76, R82 \\
Injury, poisoning and other external consequences & S10, S12, S20$^{\ast}$, S22, S30, S32, S40, S42, S52, S60, S62, S70, S72, S80, S82, S90, S91, S92, T36, T37, T39, T78$^{\ast}$, T80, T83, T84, T85 \\
External causes & X58$^{\ast}$, Y92$^{\ast}$, Y93, Y99 \\
Factors influencing health status and contact with health services & Z16, Z22, Z23, Z28, Z43, Z48, Z59, Z63, Z74, Z76, Z77, Z80, Z81, Z88, Z90, Z93, \textbf{Z96} \\
\hline
\end{tabular}
\end{table}

\begin{table}[t]
\centering
\small
\caption{No-route candidates removed by the separation screen of Appendix~\ref{sec:wm-app-anchors} and treated as codes with a route.}
\label{tab:wm-app-screen}
\begin{tabular}{p{0.93\linewidth}}
\hline
B00, B02, D18, D32, H90, J38, J39, K02, L76, N99, S37, T17, Z00, Z53, Z78 \\
\hline
\end{tabular}
\end{table}

\begin{table}[t]
\centering
\small
\caption{Rotation groups of the anchor codes outside the validation subset, stratified by the number of onsets on pairs without a cardiac change. In each rotation one group is withheld from the context estimator and given code channels.}
\label{tab:wm-app-groups}
\begin{tabular}{p{0.08\linewidth}p{0.85\linewidth}}
\hline
\textbf{Group} & \textbf{Codes} \\
\hline
1 & B37, B96, D73, E29, G57, G58, G89, H49, H91, K04, K08, K20, K27, K44, K83, L02, L30, L53, L59, M17, M21, M85, M86, N20, N83, R14, R21, R30, S20, S30, S40, S60, S82, T36, Y93, Z22, Z90 \\
2 & B97, E53, E80, G44, H04, H10, H26, J20, K50, K58, K94, L27, L89, M12, M14, M40, M41, M50, M79, N10, N32, N39, N50, N92, R52, S12, S32, S72, S80, S92, T37, T39, X58, Z16, Z23, Z77 \\
3 & B95, D12, E04, G51, H02, J30, J32, K13, K30, K43, K51, K62, K63, K82, L25, L50, L73, M51, M75, M81, M84, M90, N30, N31, R39, R49, R51, R61, R76, T78, T80, T85, Y92, Z43, Z74, Z80 \\
4 & B34, B35, B49, D25, D89, G56, H53, H57, J02, K25, K40, K42, K66, K68, K81, L29, L40, M00, M11, M19, M25, M43, M87, R11, R32, R33, R82, S10, S22, S42, S52, S62, S70, T84, Z48, Z59 \\
5 & B25, D80, E55, E74, G60, H54, H93, J01, J34, J40, K12, K26, K91, L03, L98, M15, M16, M46, M47, M48, M54, N12, N52, R10, R15, R22, R59, S90, S91, T83, Y99, Z28, Z63, Z76, Z81, Z88, Z93 \\
\hline
\end{tabular}
\end{table}
